% !TEX root = ../algorithms.tex

\section{Теорема о двойственности в линейном программировании (без доказательства леммы Фаркаша). Системы линейных неравенств. Полиномиальные алгоритмы решения лин. прог. Общая идея алгоритма внутренней точки. Преобразование задачи для нахождения какой-то внутренней точки.}

\subsection{Теорема о двойственности}

Для $A \in \RR^{m \times n}$, $b \in \RR^m$, $c \in \RR^n$ рассмотрим прямую и двойственную задачи:
\[
\begin{array}{c|c}
	\text{Прямая задача } (\mathcal{P}) & \text{Двойственная задача } (\mathcal{D}) \\ \hline
	\max\; c^T x & \min\; b^T y \\
	Ax \le b & A^T y = c \\
	& y \ge 0
\end{array}
\]

\begin{Proposition}{Слабая двойственность}{}
	Для любых допустимых $x$ (задачи $\mathcal{P}$) и $y$ (задачи $\mathcal{D}$) выполнено $c^T x \le b^T y$.
\end{Proposition}

\begin{proof}
	Так как $A^T y = c$, имеем $c^T x = (A^T y)^T x = y^T (Ax)$. Далее, из $Ax \le b$ и $y \ge 0$ следует $y^T (Ax) \le y^T b = b^T y$. Следовательно, $c^T x \le b^T y$.
\end{proof}

Любой вектор $y$ с условиями $A^T y = c$, $y \ge 0$ задаёт верхнюю оценку $b^T y$ на значение прямой задачи. Естественно искать наименьшую из таких оценок --- это и приводит к задаче $\mathcal{D}$; симметрично, двойственная к $\mathcal{D}$ снова есть $\mathcal{P}$.

\begin{Proposition}{$\mathcal{P}$ и $\mathcal{D}$ образуют двойственную пару}{}
	Задача $\mathcal{D}$ двойственна к $\mathcal{P}$, а задача $\mathcal{P}$ двойственна к $\mathcal{D}$.
\end{Proposition}

\begin{Theorem}{Сильная двойственность (теорема о двойственности)}{}
	Если одна из задач $(\mathcal{P})$ или $(\mathcal{D})$ имеет оптимальное решение, то и другая имеет оптимальное решение, причём
	\[
	c^T x^* = b^T y^*,
	\]
	где $x^*$ и $y^*$ --- оптимальные решения прямой и двойственной задач.
	
	\medskip
	Если целевая функция одной из задач не ограничена, то другая задача не имеет допустимых решений.
\end{Theorem}

\begin{proof}
	\textbf{1. Слабая двойственность} уже доказана: $c^T x \le b^T y$ для любых допустимых $x$, $y$. Отсюда, если обе задачи имеют допустимые решения, то обе ограничены, а значит, обе имеют оптимальные решения $x^*$ и $y^*$.
	
	\textbf{2. Построение оптимального $y^*$ по $x^*$.}
	
	Пусть $x^*$ --- оптимальное решение $(\mathcal{P})$. Перестановкой строк системы $Ax \le b$ разобьём ограничения на активные и неактивные:
	\[
	A = \begin{pmatrix} A' \\ A'' \end{pmatrix},
	\qquad
	b = \begin{pmatrix} b' \\ b'' \end{pmatrix},
	\quad
	A'x^* = b',\quad A''x^* < b''.
	\]
	
	Покажем, что $\nexists\, z \in \RR^n \colon A'z \le 0$ и $c^T z > 0$. В самом деле, если бы такой $z$ существовал, то для точки $x^{**} = x^* + \varepsilon z$ выполнялось бы $A'x^{**} = A'x^* + \varepsilon A'z \le b'$ (так как $A'x^* = b'$ и $A'z \le 0$), а из строгого неравенства $A''x^* < b''$ при достаточно малом $\varepsilon > 0$ следовало бы $A''x^{**} \le b''$. Значит, $Ax^{**} \le b$, то есть $x^{**}$ допустима, но $c^T x^{**} = c^T x^* + \varepsilon c^T z > c^T x^*$ --- противоречие с оптимальностью $x^*$.
	
	Следовательно, для матрицы $A'$ второе условие леммы Фаркаша не выполнено, поэтому выполнено первое:
	\[
	\exists\, y' \ge 0\colon (A')^T y' = c.
	\]
	Положим $y = \begin{pmatrix} y' \\ 0 \end{pmatrix} \in \RR^m$. Тогда $y \ge 0$ и $A^T y = (A')^T y' + (A'')^T 0 = c$, то есть $y$ --- допустимое решение $(\mathcal{D})$.
	
	Проверим значения целевых функций:
	\[
	b^T y = (b')^T y' = (A'x^*)^T y' = (x^*)^T (A')^T y' = (x^*)^T c = c^T x^*.
	\]
	По слабой двойственности $b^T y = c^T x^*$ --- наибольшее достижимое значение, поэтому $y$ оптимально в $(\mathcal{D})$, а $x^*$ --- в $(\mathcal{P})$.
	
	\textbf{3. Симметричность.}
	
	Покажем, что если оптимальное решение имеет задача $(\mathcal{D})$, то и $(\mathcal{P})$ тоже. Перепишем $(\mathcal{D})$ в виде $\max\,(-b^T y)$, а ограничения $A^T y = c$, $y \ge 0$ --- системой неравенств:
	\[
	\max\; (-b^T y)
	\quad \text{при} \quad
	\begin{cases}
		A^T y \le c,\\
		-A^T y \le -c,\\
		-y \le 0.
	\end{cases}
	\]
	Задача, двойственная к этой, в точности эквивалентна исходной задаче $(\mathcal{P})$: $\max\; c^T x$, $Ax \le b$. Применяя к ней ту же конструкцию, по оптимальному решению $(\mathcal{D})$ строим допустимое (и потому оптимальное) решение $(\mathcal{P})$.
\end{proof}

\subsection{Системы линейных неравенств}

\begin{Remark}{Сведение ЛП к задаче выполнимости}{}
	Пусть имеется алгоритм, за полиномиальное время решающий задачу выполнимости системы линейных неравенств (находит решение либо сообщает, что решений нет). Тогда нахождение оптимального решения задачи ЛП сводится к проверке выполнимости одной системы. Действительно, пары $(x, y)$ являются оптимальными решениями $(\mathcal{P})$ и $(\mathcal{D})$ тогда и только тогда, когда выполнена система
	\[
	Ax \le b,\qquad A^T y = c,\qquad y \ge 0,\qquad b^T y - c^T x \le 0.
	\]
	По слабой двойственности для любых допустимых $x$, $y$ всегда $b^T y - c^T x \ge 0$, поэтому добавленное неравенство $b^T y - c^T x \le 0$ эквивалентно равенству $b^T y = c^T x$, то есть условию оптимальности.
	
	При желании равенство $A^T y = c$ заменяется парой неравенств $A^T y \le c$ и $-A^T y \le -c$. Таким образом, задача ЛП полностью сводится к проверке выполнимости системы линейных неравенств, и наоборот.
\end{Remark}

\subsection{Полиномиальные алгоритмы}

\begin{Definition}{Длина входа}{}
	Пусть все элементы $A$, $b$, $c$ рациональны и кодируются не более чем $w$ битами каждый. Тогда длиной входа называется
	\[
	L = (mn + m + n)\, w = \Theta(mnw).
	\]
\end{Definition}

\begin{Definition}{Полиномиальный алгоритм}{}
	Алгоритм решения задачи ЛП называется \emph{полиномиальным}, если его время работы равно $O(L^{O(1)})$.
\end{Definition}

Симплекс-метод обходит вершины многогранника, двигаясь по рёбрам в направлении вектора $c$. Большинство его вариантов имеют экспоненциальную сложность в худшем случае, но хорошо работают на практике. Полиномиальные алгоритмы: метод эллипсоидов (Хачиян, 1979) и метод внутренней точки (Кармаркар, 1984).

\subsection{Метод внутренней точки}

\begin{Definition}{Внутренняя точка}{}
	Точка $x_0$ называется \emph{внутренней допустимой точкой} для системы $Ax \le b$, если $Ax_0 < b$ покомпонентно.
\end{Definition}

\subsubsection*{Нахождение внутренней точки}

\begin{Remark}{Как найти внутреннюю точку}{}
	Для поиска внутренней точки исходной системы $Ax \le b$ рассмотрим вспомогательную задачу
	\[
	\max \; t
	\quad \text{при} \quad
	Ax + t\,\mathbf{1} \le b,
	\]
	где $\mathbf{1} = (1,\dots,1)^T \in \RR^m$. Если её оптимальное значение $t^* > 0$, то соответствующая точка $x^*$ удовлетворяет $Ax^* < b$, то есть является внутренней для исходной системы.
	
	У вспомогательной задачи всегда есть очевидная внутренняя точка, например
	\[
	x = 0,\qquad t = -\max_{1 \le i \le m} |b_i| - 1,
	\]
	поскольку тогда $t < b_i$ для всех $i$, а значит $Ax + t\,\mathbf{1} = t\,\mathbf{1} < b$.
\end{Remark}

\subsubsection*{Барьерная задача и центральный путь}

Пусть строки матрицы $A$ имеют вид $a_1^T,\dots,a_m^T$, а $s_i(x) := b_i - a_i^T x$ обозначают зазоры до ограничений.

\begin{Definition}{Логарифмический барьер}{}
	Для $x$ во внутренности многогранника ($Ax < b$) положим
	\[
	g(x) := \sum_{i=1}^m \ln s_i(x) = \sum_{i=1}^m \ln(b_i - a_i^T x).
	\]
	Для параметра $w > 0$ рассмотрим \textbf{барьерную задачу}
	\[
	\max\; \bigl(w\, c^T x + g(x)\bigr)
	\quad \text{при} \quad
	Ax \le b.
	\]
	Одно из её решений обозначим $x_w^*$.
\end{Definition}

Смысл логарифмических слагаемых: при приближении $x$ к границе многогранника какой-нибудь зазор $s_i(x) \to 0^+$, поэтому $g(x) \to -\infty$. Барьер $g$ «отталкивает» от границы, и максимизатор $x_w^*$ остаётся строго внутри допустимой области. При этом слагаемое $w\,c^T x$ тянет точку в сторону роста целевой функции; баланс между ними и определяет положение $x_w^*$.

\begin{Proposition}{Центральный путь}{}
	Семейство $\{x_w^*\}_{w > 0}$ образует \emph{центральный путь}. Поскольку
	\[
	\arg\max_{Ax < b}\bigl(w\,c^T x + g(x)\bigr) = \arg\max_{Ax < b}\!\left(c^T x + \tfrac{1}{w}\,g(x)\right),
	\]
	множитель при барьере есть $\tfrac{1}{w}$. При $w \to +\infty$ этот множитель убывает, влияние барьера ослабевает, и точки $x_w^*$ приближаются к оптимальному решению исходной задачи; наоборот, при $w \to 0^+$ доминирует барьер, и $x_w^*$ стремится к «аналитическому центру» многогранника.
\end{Proposition}

\begin{figure}[H]\centering
	\begin{tikzpicture}[>=Stealth, font=\small]
		% многогранник P
		\coordinate (P1) at (0,0);
		\coordinate (P2) at (4,0.3);
		\coordinate (P3) at (4.6,2.2);
		\coordinate (P4) at (2.4,3.6);
		\coordinate (P5) at (-0.4,2.0);
		\fill[conecolor] (P1)--(P2)--(P3)--(P4)--(P5)--cycle;
		\draw[axiscolor, thick] (P1)--(P2)--(P3)--(P4)--(P5)--cycle;
		\node[axiscolor] at (1.7,0.85) {$P=\{Ax\le b\}$};
		
		% направление c
		\draw[thickvec, color=cvec] (3.0,3.0) -- ++(1.1,0.33)
		node[above right, color=cvec] {$c$};
		
		% аналитический центр (w -> 0)
		\coordinate (AC) at (1.8,1.75);
		\filldraw[dvec] (AC) circle (1.6pt);
		\node[dvec, below left] at (AC) {аналит. центр ($w\to 0$)};
		
		% центральный путь
		\draw[dvec, thick] (AC) .. controls (3.1,1.5) and (4.0,1.78) .. (4.45,2.10);
		\filldraw[dvec] (3.25,1.59) circle (1.4pt);
		\filldraw[dvec] (4.05,1.85) circle (1.4pt);
		\node[dvec] at (3.45,2.65) {центральный путь $\{x_w^*\}$};
		
		% оптимальная вершина
		\filldraw[cvec] (P3) circle (2pt);
		\node[cvec, right] at (P3) {$x^*$ ($w\to\infty$)};
	\end{tikzpicture}
	\caption{Центральный путь $\{x_w^*\}$ внутри многогранника $P$. При малом $w$ доминирует барьер и точка близка к аналитическому центру; при $w \to \infty$ вес барьера $\tfrac{1}{w}$ убывает, и центральные точки движутся вдоль пути к оптимальной вершине $x^*$. Метод следует именно вдоль этой кривой.}
\end{figure}

\subsubsection*{Идея алгоритма}

Метод следует вдоль центрального пути. Начав с внутренней точки $x_0$, он постепенно увеличивает параметр $w$ (то есть уменьшает вес барьера $\tfrac{1}{w}$) и на каждом шаге сдвигает текущую точку ближе к соответствующей центральной точке $x_w^*$. Поскольку при $w \to +\infty$ центральные точки приближаются к оптимуму исходной задачи, итерации сходятся к оптимальному решению. При аккуратном выборе шага метод даёт полиномиальный алгоритм решения задачи ЛП.